% ===========================================================================
% Supplement S6 — presentation methods for the validated-scenario video
% artifacts: the C² log-domain time remap and shutter synthesis. Sources
% of truth: the shipped implementation (web/app/framestep.js —
% smootherstep / logRemapFrames / shutterSubframes, the one frame-timeline
% module shared by the in-app exporter and the node render harness) and
% the cited primary sources (entry-level metadata verified against them on
% 2026-07-26; see references.bib).
% SCOPE GUARD (taxonomy discipline): the simulator's
% output of record is the kernel event stream (Section 3 of the main
% text); this section documents how gated records are PRESENTED as video,
% and nothing in it enters the paper's V&V taxonomy or evidence chain.
% ===========================================================================
\section{Presentation methods for the validated-scenario video artifacts}
\label{sec:supp-timeremap}

The artifact bundle's validated-scenario video pack
(Section~\ref{sec:appendix-repro}) renders sixteen catalog entries ---
each a solved strike, or a short ladder of strikes, replayed through the
canonical binary shipped at the time of its render --- as annotated
slow-motion videos. The simulator's output of record is the
kernel event stream, not any sampled rendering of it (main text,
Section~3), and no claim of the paper rests on these videos: each entry
first re-simulates its claimed scenario on that shipped canonical binary
and must reproduce its pinned record digest \emph{before} any frame is
captured, so the videos are presentations of already-verified records.
The per-entry kernel stamp is recorded in the entry's gate report; the
pack's stamps postdate the paper's battery pin along the supersession
lineage of the main text, and each entry's digests are pinned against
its own stamp.
This section documents the deterministic time-remap and shutter
synthesis used by those presentations, for two reasons: the frame
schedule is deterministic --- a pure function of the entry's playback
plan, which is authored in the render tool's scenario catalog, with the
derived schedule's summary (frame counts, worst rate-step ratio,
shutter parameters, replay spans) pinned in the per-entry gate
report --- and the design choices are
non-obvious enough that a reproducer would otherwise have to
reverse-engineer them from the delivered bytes.

\paragraph{Frame schedule and the log-domain speed ramp.}
Videos are delivered at a constant 60 frames per second; one output
frame advances the kernel clock by $\Delta = 1/60$~s times a
dimensionless playback rate $r$ (kernel seconds consumed per output
second; $r = 1$ is real time). A playback plan is a short list of rate
segments over kernel time, each pair of consecutive segments separated
either by a camera cut (a hard rate step on the cut frame) or by a
\emph{ramp}. A hard rate step without a cut reads as a speed jerk, so
cut-less boundaries between rates $r_1$ and $r_2$ are ramped in the
logarithmic rate domain:
\begin{equation}
r(u) \;=\; r_1\,\bigl(r_2/r_1\bigr)^{S(u)},
\qquad
S(x) \;=\; x^3\,(6x^2 - 15x + 10),
\qquad u \in [0,1],
\label{eq:supp-lograte}
\end{equation}
where $S$ is the quintic ``smootherstep'' polynomial, with zero first
and second derivatives at both endpoints; the remap is therefore $C^2$
in log-rate, so the on-screen speed carries no visible acceleration
kink. $S$ is, term for term, the minimum-jerk rest-to-rest trajectory of
human reaching movements \citep[eq.~(2), p.~1689]{flash1985coordination}:
shaping the ramp with it makes the log-rate trajectory the minimum-jerk
path between $\log_2 r_1$ and $\log_2 r_2$. The ramp duration follows a
per-octave law,
\begin{equation}
D \;=\; \operatorname{clamp}\bigl(0.30 \cdot \lvert \log_2 (r_2/r_1)
\rvert,\; 0.5,\; 1.5\bigr)~\text{output seconds},
\label{eq:supp-octave}
\end{equation}
quantized to $n = \max(1, \operatorname{round}(D/\Delta))$ whole output
frames --- a fixed time cost per doubling of speed, bounded below and
above so that small steps are not abrupt and large ones do not drag.
The per-frame kernel advance integrates the ramp at the frame midpoint,
$\mathrm{d}t_j = \Delta \cdot r((j + \tfrac12)/n)$ for
$j = 0,\dots,n-1$, and the ramp is scheduled as a \emph{lead-in}: it
starts on the last steady frame from which its integrated kernel
consumption still completes at or before the segment boundary, so that
whenever the outgoing segment can host the ramp --- the common case ---
the target rate is established at or before the boundary and impacts
land fully inside their slow-motion rate. A segment too short to host
its ramp instead starts the ramp at segment entry and overshoots the
boundary; the published pack contains one such replay pass, and the
exception is why the rate-step law below and the per-entry gate windows
are measured on the \emph{emitted} clock, never on the nominal plan.
Plan validation rejects
non-positive or non-finite rates, which the log-domain form requires
(the analogous restriction in the Web Audio specification is discussed
below).

\paragraph{Perceptual basis for ramping in the log domain.}
The ramp is built in the logarithmic speed domain, where equal steps are
equal speed \emph{ratios}; the quintic easing modulates how fast that
axis is traversed (its rate reaching zero at both endpoints), so the
constant-ratio character lives in the easing variable $S(u)$ rather
than literally in each output frame. The ratio axis is the perceptually
relevant one for visual speed. The classical measurement is Weber-law
behavior: the
just-noticeable difference in speed grows in proportion to the reference
speed \citep{brown1960weber}. \citet{debruyn1988velocity} measured the
Weber fraction for velocity discrimination at approximately 7\% at the
optimal speeds (84\%-correct staircase criterion), roughly constant over
a 4--64~deg/s plateau and rising (U-shaped) at slower and faster
speeds. In modeling terms, \citet{stocker2006noise} adopt a logarithmic
speed representation for their likelihood and internal-noise functions
(a log-Normal likelihood, $\tilde{v} = \ln((v + v_0)/v_0)$, in their
Methods), motivated by exactly this Weber--Fechner character of speed
discrimination --- a modeling choice for the noise function, not a
direct neural-coding proof; the neural grounding of a logarithmic,
scale-invariant speed representation is the middle-temporal-area
population analysis of \citet{nover2005logarithmic}. A perceptually
equal speed change is therefore an equal step in $\log r$, and
Eq.~\eqref{eq:supp-lograte} ramps along exactly that axis, its
minimum-jerk easing distributing the steps in time. The same reasoning is
codified for audio in the W3C Web Audio API, whose
\code{exponentialRampToValueAtTime} automation has the identical
functional form $v(t) = V_0 (V_1/V_0)^{(t - T_0)/(T_1 - T_0)}$ and whose
specification recommends exponential ramps for ``filter frequencies and
playback rate . . . because of the way humans perceive sound''
\citep{webaudio_spec}. The specification's domain restriction is
instructive: an exponential ramp can neither reach nor pass through
zero --- the W3C text throws a \code{RangeError} only when the target
value equals zero, while a zero or opposite-signed \emph{start} value
silently freezes the ramp (a ramp between two negative values is
permitted). Our plan validator is intentionally stricter than that
domain: playback rates must be finite and strictly positive, so the
degenerate cases the specification handles by freezing or throwing
cannot be expressed in a plan at all.

\paragraph{Shutter synthesis over world-time windows.}
Motion blur is synthesized, not post-processed: each delivered frame is
the accumulation of 8 sub-frame renders spanning a fixed
$310^\circ$ shutter angle, i.e.\ the shutter opens at the frame's kernel
time $t_j$ and stays open for $(310/360) \cdot \mathrm{d}t_j$ of
\emph{kernel} (world) time, where $\mathrm{d}t_j$ is the integrated
per-frame advance above --- never the nominal segment rate. Blur length
therefore follows the on-screen speed through ramps exactly as a
physical camera filming the slow-motion playback would behave. In the
frame domain the sub-frame animation steps sum to exactly $\Delta$ per
output frame (the first sub-frame carries the closed-shutter
remainder), so frame counts and presentation-layer animation lengths
are independent of the shutter setting.

\paragraph{Monotone-map anchoring.}
The schedule is emitted in monotone runs: within a run the map from
output frame index to kernel time is strictly increasing, and a
broadcast-style replay (a rewind to re-cover an impact) starts a new
run, its boundary marked by a camera cut and a 0.4~s cross-dissolve
identifying re-covered time. The \emph{placement} of timed presentation
content ---
predicted-overlay geometry, measured-value callouts, story-beat text
reveals --- is declared as \emph{kernel-time} windows and resolved to
output frames through this monotone map, never positioned by hand in
output seconds; within a window, the reveal choreography itself (alpha
fades, type entrances) plays authored fixed output-time durations ---
0.4~s fades matching the dissolve above --- and a few entries carry
small authored output-time leads. An annotation window that would match
more than one run of its
camera is rejected rather than guessed at. One mechanical law bounds
the emitted clock: the worst consecutive cut-less per-frame
kernel-advance ratio, measured on the emitted clock per monotone run,
must not exceed $1.35\times$; the measured value is recorded in the
per-entry gate report, and an entry that exceeds the bound fails before
any frame is captured.

\paragraph{Relation to prior art.}
The two ingredients of Eq.~\eqref{eq:supp-lograte} appear separately in
published work. Speed-template curves in video-editing tools enumerate
logarithmic and exponential playback-speed transitions among other
shapes, motivated by editorial smoothness rather than psychophysics
\citep{gopro2017variable,adobe2012changing}; and Weber--Fechner-guided
playback-rate control appears in adaptive fast-forward summarization,
where perceived fast-forward speed is modeled on a logarithmic axis
\citep{jin2017elasticplay}. We have not found a published description
that combines a perceptually motivated log-domain speed curve with
$C^2$ minimum-jerk easing for cinematic slow-motion time remapping; we
record the synthesis here as presentation methodology --- it is not a
claim of this paper, and none of the paper's results depend on it.
